\chapter{Conclusion}
\label{ch:conclusion}

% ============================================================================
% This chapter concludes the thesis and outlines future work
% ============================================================================

\section{Summary}
\label{sec:conclusion_summary}

This thesis presented a complete sleep stage classification pipeline using machine learning models trained on EEG data from the BOAS dataset. The main contributions are:

\begin{enumerate}
    \item \textbf{Complete pipeline implementation:} A modular Python system (25 modules, 16,571 lines of code) for sleep stage classification from raw EEG data to predictions.

    \item \textbf{Two-tier caching system:} A fingerprint-based caching architecture with feature-level and model-level caches that significantly reduces computational overhead while ensuring reproducibility.

    \item \textbf{LOSO cross-validation:} Subject-independent evaluation on 128 subjects using Leave-One-Subject-Out cross-validation, ensuring honest generalization assessment.

    \item \textbf{Model comparison:} Empirical comparison of XGBoost and Random Forest across different feature selection configurations (top-K = 30, 50, 75).

    \item \textbf{Open-source implementation:} Complete, documented code for reproducible research.
\end{enumerate}

\section{Key Findings}
\label{sec:key_findings}

% TODO: Fill in after experiments

The experimental results demonstrated:

\begin{enumerate}
    \item [Model X] achieved [Y\%] accuracy with Cohen's kappa of [Z], demonstrating [good/moderate] subject-independent performance.

    \item [Feature selection finding - optimal top-K, most important features]

    \item The two-tier caching system reduced total experiment time by [X\%], enabling rapid iteration during development.

    \item [Other key finding - e.g., which features are most important, which stages are hardest to classify]
\end{enumerate}

\section{Contributions to the Field}
\label{sec:contributions_field}

This work contributes to sleep stage classification research by:

\begin{itemize}
    \item Providing a reproducible baseline for feature-based machine learning on the BOAS dataset
    \item Demonstrating the value of rigorous LOSO cross-validation for subject-independent evaluation
    \item Introducing a practical caching system architecture for computationally expensive ML experiments
    \item Offering open-source code for future research
\end{itemize}

\section{Lessons Learned}
\label{sec:lessons_learned}

Key lessons from this project:

\begin{enumerate}
    \item \textbf{Caching is critical:} For LOSO cross-validation with many subjects, intelligent caching transforms experiment turnaround from days to minutes.

    \item \textbf{LOSO is stringent:} Leave-One-Subject-Out cross-validation reveals true subject-independent performance, which is typically lower than random K-fold results.

    \item \textbf{Feature selection matters:} ANOVA-based selection is fast and effective; 200× faster than mutual information with <1\% accuracy loss.

    \item \textbf{Class imbalance is real:} Rare stages like N1 are consistently harder to classify, affecting macro-averaged metrics.

    \item \textbf{Reproducibility requires discipline:} Fingerprint-based cache invalidation, fixed random seeds, and comprehensive logging are essential.
\end{enumerate}

\section{Limitations}
\label{sec:conclusion_limitations}

The main limitations of this work are:

\begin{enumerate}
    \item \textbf{Single dataset:} Evaluation on BOAS only; generalization to other datasets is uncertain.

    \item \textbf{Healthy subjects:} No validation on clinical populations with sleep disorders.

    \item \textbf{Handcrafted features:} Feature engineering requires domain knowledge; deep learning might discover better representations automatically.

    \item \textbf{Limited hyperparameter search:} Computational constraints limited hyperparameter optimization.

    \item \textbf{No temporal context:} Single-epoch classification ignores sequential patterns across neighboring epochs.
\end{enumerate}

\section{Future Work}
\label{sec:future_work}

Several directions for future research:

\subsection{Model Improvements}
\label{subsec:future_models}

\begin{itemize}
    \item \textbf{Deep learning:} Implement CNNs or RNNs for automatic feature extraction and temporal modeling
    \item \textbf{Sequence models:} Use LSTMs or Transformers to model temporal context across epochs
    \item \textbf{Class imbalance handling:} Apply SMOTE, class weights, or focal loss to improve minority class performance
    \item \textbf{Hyperparameter optimization:} Use Bayesian optimization or grid search with the caching system
\end{itemize}

\subsection{Dataset Extensions}
\label{subsec:future_datasets}

\begin{itemize}
    \item \textbf{Multi-dataset validation:} Evaluate on Sleep-EDF, MASS, SHHS to assess generalization
    \item \textbf{Clinical populations:} Test on sleep disorder patients (sleep apnea, insomnia, etc.)
    \item \textbf{Cross-dataset transfer:} Train on one dataset, test on another (domain adaptation)
    \item \textbf{Multi-modal fusion:} Combine EEG with other signals (EOG, EMG, ECG)
\end{itemize}

\subsection{System Enhancements}
\label{subsec:future_system}

\begin{itemize}
    \item \textbf{GPU acceleration:} Implement XGBoost GPU support (\texttt{tree\_method='gpu\_hist'}) for faster training
    \item \textbf{Distributed computing:} Parallelize LOSO folds across machines for massive speedup
    \item \textbf{Real-time classification:} Optimize for online sleep staging during recording
    \item \textbf{Interactive visualization:} Build web dashboard for exploring results
\end{itemize}

\subsection{Research Questions}
\label{subsec:future_research}

\begin{itemize}
    \item \textbf{Feature importance stability:} Do selected features generalize across datasets?
    \item \textbf{Subject clustering:} Can we identify subject groups with similar sleep patterns?
    \item \textbf{Interpretability:} Use SHAP or LIME to explain model predictions
    \item \textbf{Uncertainty quantification:} Estimate prediction confidence for ambiguous epochs
\end{itemize}

\section{Closing Remarks}
\label{sec:closing_remarks}

Sleep stage classification is a fundamental task in sleep medicine and neuroscience research. This thesis demonstrated that feature-based machine learning with rigorous LOSO cross-validation can achieve [good/competitive] subject-independent performance on the BOAS dataset.

The two-tier caching system introduced here addresses a practical challenge in machine learning research: balancing computational efficiency with reproducibility. By using cryptographic fingerprints to uniquely identify configurations, the system ensures that experiments are recomputed when parameters change while reusing results when configurations are identical.

As sleep monitoring transitions from clinical labs to consumer devices, automated sleep stage classification will become increasingly important. Future work should focus on validating these methods on diverse populations, improving temporal modeling, and adapting pipelines for real-time deployment.

The complete implementation is available as open-source code, enabling future researchers to build on this foundation for sleep stage classification and beyond.

% ============================================================================
% Final word count estimate: ~800 words (this chapter)
% Total thesis estimate: ~15,000-20,000 words
% ============================================================================
